\subsection{谱测度}

在本节中，我们固定一个复希尔伯特空间 E。

回顾：若 A 是交换含幺星代数，我们用 X(A) 表示其幺特征标构成的紧拓扑空间（I, p. 29 的推论 1），并用 \(\mathcal{G}_{\mathrm{A}}\) 表示 A 的盖尔范德变换（I, p. 7 的定义 5），它是从 A 到 \(\mathcal{C}(\mathrm{X}(\mathrm{A}))\) 上的含幺对合代数的等距同构（I, p. 108 的定理 1）。我们将用 \(\mathcal{H}_{\mathrm{A}}\) 表示其逆同构。

引理 9. --- 设 A 是 \(\mathcal{L}(\mathrm{E})\) 的交换含幺星子代数。对 E 中所有 \(x\) 与 \(y\) ，映射 \(\mu\) ——从 \(\mathsf{X}(\mathrm{A})\) 到 \(\mathbf{C}\) 中，由 \(\mu(f) = \langle x | \mathcal{H}_{\mathrm{A}}(f)y \rangle\) （对每个 \(f \in \mathcal{C}(\mathsf{X}(\mathrm{A}))\) ）所定义——是紧空间 \(\mathsf{X}(\mathrm{A})\) 上的有界测度。当 \(x = y\) 时它是正的。其范数 \(\leqslant \|x\|\ \|y\|\) ，且当 \(x = y\) 时取等号。

映射 \(\mu\) 是线性的。对每个函数 \(f\in\mathcal{C}(\mathsf{X}(\mathsf{A}))\) ，我们有

\[
| \mu (f) | \leqslant \| x \| \| y \| \| \mathcal {H} _ {\mathrm{A}} (f) \| = \| x \| \| y \| \| f \|,
\]

因而 \(\mu\) 是 \(\mathsf{X}(\mathsf{A})\) 上范数 \(\leqslant \| x\|\| y\|\) 的有界测度。

假设 \(x = y\) 。对每个正函数 \(f \in \mathcal{C}(\mathsf{X}(\mathsf{A}))\) ，元素 \(\mathcal{H}_{\mathsf{A}}(f)\) 是 A 的正元素（I, p. 115 的例 1），因而是 \(\mathcal{L}(\mathsf{E})\) 的正元素，于是 \(\mu(f) = \langle x | \mathcal{H}_{\mathsf{A}}(f)(x) \rangle \geqslant 0\) （I, p. 138 的命题 8）。这表明 \(\mu\) 是正测度。由于 \(\mu(1) = \|x\|^2\) ，\(\mu\) 的总质量是 \(\|x\|^2\) （INT, III, p. 58, § 1, no. 8，推论 2）。

定义 2. --- 设 A 是 \(\mathcal{L}(\mathrm{E})\) 的交换含幺星子代数。对 E 中所有 \(x\) 与 \(y\) ，线性形式 \(f \mapsto \langle x | \mathcal{H}_{\mathrm{A}}(f)y \rangle\) （\(\mathcal{C}(\mathrm{X}(\mathrm{A}))\) 上的）称为 \((x,y)\) 关于 A 的谱测度。若 \(x = y\) ，我们称它为 \(x\) 关于 A 的谱测度。

注. --- 设 A 是 \(\mathcal{L}(\mathrm{E})\) 的交换含幺星子代数。对 E 中所有 \(x\) 与 \(y\) ，用 \(\mu_{x,y}\) 表示 \((x,y)\) 关于 A 的谱测度。由 \((x,y) \mapsto \mu_{x,y}\) 定义的从 \(\mathrm{E} \times \mathrm{E}\) 到 \(\mathcal{M}^1(\mathrm{X}(\mathrm{A}))\) 中的映射是半双线性的。

设 \(u\) 是 E 的正规自同态，A 是 \(\mathcal{L}(\mathrm{E})\) 的、由 \(u\) 所生成的含幺星子代数。它是交换的。空间 \(\mathsf{X}(\mathrm{A})\) 与 \(\mathrm{Sp}(u)\) 等同——通过映射 \(\chi \mapsto \chi(u)\) （I, p. 109 的引理 10）。于是，谱测度 \(\mu_{x,y}\) ——\((x,y)\) 关于 A 的——与 \(\mathrm{Sp}(u)\) 上的一个测度等同，该测度称为 \((x,y)\) 关于 \(u\) 的谱测度。于是，对每个函数 \(f \in \mathcal{C}(\mathrm{Sp}(u))\) ，我们有

\[
\int_ {\mathrm{Sp} (u)} f d \mu_ {x, y} = \langle x | f (u) y \rangle
\]

依据 I, p. 110 的定义 5。

